\documentclass[../../../include/open-logic-section]{subfiles}

\begin{document}

\olfileid[id]{sth}{cardinals}{cardsasords}
\olsection{Kardinal sebagai Ordinal}

Pada kenyataannya, teori kardinal kita akan memanfaatkan teori ordinal kita
(tanpa malu-malu). Artinya, kita hanya akan mendefinisikan kardinal sebagai
ordinal-ordinal tertentu. Secara khusus, kita mengajukan definisi berikut:

\begin{defn}\ollabel{defcardinalasordinal}
Jika $A$ dapat diurutkan dengan baik, maka $\card{A}$ ialah ordinal terkecil
$\gamma$ sedemikian sehingga $\cardeq{A}{\gamma}$. Untuk sebarang ordinal
$\gamma$, kita mengatakan bahwa $\gamma$ merupakan suatu \emph{kardinal}
jika dan hanya jika $\gamma = \card{\gamma}$.
\end{defn}

Kita baru saja menggunakan ungkapan ``$A$ dapat diurutkan dengan baik''.
Seperti hampir selalu terjadi dalam matematika, ungkapan modal di sini
hanyalah uraian informal bagi suatu klaim eksistensial: mengatakan ``$A$
dapat diurutkan dengan baik'' hanyalah mengatakan ``terdapat suatu relasi
yang mengurutkan $A$ dengan baik''.

Namun, terdapat masalah pada \olref{defcardinalasordinal}. Kita menginginkan
agar \emph{setiap} himpunan mempunyai ukuran, yakni agar $\card{A}$ ada untuk
setiap~$A$. Akan tetapi, definisi yang baru saja kita berikan diawali suatu
syarat: ``\emph{Jika} $A$ dapat diurutkan dengan baik\ldots''. Jika terdapat
suatu himpunan $A$ yang tidak dapat diurutkan dengan baik, definisi kita akan
gagal mendefinisikan objek~$\card{A}$.

Jadi, untuk menggunakan \olref{defcardinalasordinal}, kita memerlukan jaminan
bahwa setiap himpunan dapat diurutkan dengan baik. Sayangnya, jaminan ini
tidak tersedia dalam~$\ZF$. Maka, jika kita hendak menggunakan
\olref{defcardinalasordinal}, tidak ada pilihan selain menambahkan suatu
aksioma baru, misalnya:
\begin{axiom}[Pengurutan Baik]
Setiap himpunan dapat diurutkan dengan baik.
\end{axiom}
Kita akan membahas apakah Aksioma Pengurutan Baik dapat diterima dalam
\olref[choice][]{chap}. Namun, mulai sekarang kita akan langsung
menggunakannya. Dengan aksioma tersebut, cukup mudah untuk membuktikan bahwa
kardinal (sebagaimana didefinisikan dalam \olref{defcardinalasordinal}) ada
dan berperilaku dengan baik:

\begin{lem}\ollabel{lem:CardinalsExist}
Untuk setiap himpunan $A$:
\begin{enumerate}
	\item\ollabel{cardaexists} $\card{A}$ ada dan unik;
	\item\ollabel{cardaapprox}  $\cardeq{\card{A}}{A}$;
	\item\ollabel{cardaidem}  $\card{A}$ merupakan suatu kardinal, yakni
	$\card{A} = \card{\card{A}}$;
\end{enumerate}
\end{lem}

\begin{proof}
Tetapkan $A$. Berdasarkan Pengurutan Baik, terdapat suatu pengurutan baik
$\tuple{A, R}$. Berdasarkan
\olref[ordinals][ordtype]{thmOrdinalRepresentation}, $\tuple{A, R}$
isomorfik dengan suatu ordinal unik, $\beta$. Jadi
$\cardeq{A}{\beta}$. Berdasarkan Induksi Transfinit, terdapat suatu ordinal
terkecil yang unik, $\gamma$, sedemikian sehingga
$\cardeq{A}{\gamma}$. Jadi $\card{A} = \gamma$, yang menetapkan
\olref{cardaexists} dan \olref{cardaapprox}. Untuk menetapkan
\olref{cardaidem}, perhatikan bahwa jika $\delta \in \gamma$, maka
$\cardless{\delta}{A}$ berdasarkan pilihan kita atas $\gamma$, sehingga juga
$\cardless{\delta}{\gamma}$ karena kesamaan kardinalitas merupakan relasi
ekuivalensi (\olref[sfr][siz][equ]{equinumerosityisequi}). Jadi $\gamma =
\card{\gamma}$.
%So, for reductio, suppose that there is some ordinal $\delta \in \gamma$ such that $\cardeq{\gamma}{\delta}$. Then, $\cardeq{\cardeq{A}{\gamma}}{\delta}$ so that $\cardeq{A}{\delta}$ by \olref[sfr][set][equ]{equinumerosityisequi}, which contradicts the choice of $\gamma$ as the least ordinal such that $\cardeq{A}{\gamma}$.
\end{proof}

Hasil berikutnya menjamin Prinsip Cantor, bahkan lebih dari itu. (Perhatikan
bahwa kardinal mewarisi pengurutannya dari ordinal, yakni
$\cardfont{a} < \cardfont{b}$ jika dan hanya jika
$\cardfont{a} \in \cardfont{b}$. Dalam merumuskan hal ini, kita akan
menggunakan huruf Fraktur untuk objek yang kita ketahui sebagai kardinal.
Konvensi ini cukup baku. Alternatif yang lazim ialah menggunakan huruf Yunani,
sebab kardinal merupakan ordinal, tetapi memilihnya dari bagian tengah
alfabet, misalnya $\kappa, \lambda$.):
\begin{lem}\ollabel{lem:CardinalsBehaveRight}
Untuk sebarang himpunan $A$ dan $B$:
\begin{align*}
	\cardeq{A}{B} &\text{ jika dan hanya jika } \card{A} = \card{B}\\
	\cardle{A}{B} &\text{ jika dan hanya jika } \card{A} \leq \card{B}\\
	\cardless{A}{B}&\text{ jika dan hanya jika } \card{A} < \card{B}
\end{align*}
\end{lem}

\begin{proof}
Kita akan membuktikan arah kiri-ke-kanan dari klaim kedua (kasus-kasus lain
serupa dan diserahkan sebagai latihan). Jadi, perhatikan diagram berikut:
\begin{center}
	\begin{tikzpicture}
	\node (nodea) {$A$};
	\node[right = 6em of nodea] (nodeb) {$B$};
	\node[below = 2em of nodea] (nodecarda) {$\card{A}$};
	\node[below = 2em of nodeb] (nodecardb) {$\card{B}$};
	\draw[->] (nodea)--(nodeb);
	\draw[<->] (nodea)--(nodecarda);
	\draw[<->] (nodeb)--(nodecardb);
	\draw[->, dashed] (nodecarda)--(nodecardb);
\end{tikzpicture}
\end{center}
Panah berkepala dua menunjukkan !!{bijection}s, yang keberadaannya dijamin
oleh \olref{lem:CardinalsExist}. Dengan mengandaikan
$\cardle{A}{B}$, terdapat !!a{injection} $A\to B$. Sekarang, dengan
mengikuti panah dari $\card{A}$ ke $A$, lalu ke $B$, lalu ke $\card{B}$, kita
memperoleh !!a{injection} $\card{A} \to \card{B}$ (panah putus-putus).
\end{proof}\noindent Kita juga dapat menggunakan
\olref{lem:CardinalsBehaveRight} untuk membuktikan kembali
Schr\"{o}der--Bernstein. Klaim ini menyatakan bahwa jika
$\cardle{A}{B}$ dan $\cardle{B}{A}$, maka $\cardeq{A}{B}$. Kita menyatakannya
sebagai \olref[sfr][siz][sb]{thm:schroder-bernstein}, tetapi pertama kali
membuktikannya---dengan usaha yang cukup besar---dalam
\olref[sfr][infinite][card-sb]{sec}. Sekarang perhatikan:

\begin{proof}[Pembuktian kembali Schr\"oder--Bernstein]
Jika $\cardle{A}{B}$ dan $\cardle{B}{A}$, maka $\card{A} \leq \card{B}$
dan $\card{B} \leq \card{A}$ berdasarkan
\olref{lem:CardinalsBehaveRight}. Jadi $\card{A} = \card{B}$ dan
$\cardeq{A}{B}$ berdasarkan Trikotomi dan
\olref{lem:CardinalsBehaveRight}.
\end{proof}
\noindent
Walaupun bukti ini sangat sederhana, bukti tersebut secara implisit
bergantung baik pada Penggantian (untuk memperoleh
\olref[ordinals][ordtype]{thmOrdinalRepresentation}) maupun pada Pengurutan
Baik (untuk menjamin \olref{lem:CardinalsBehaveRight}). Sebaliknya, bukti
dalam \olref[sfr][infinite][card-sb]{sec} jauh lebih mandiri (bahkan dapat
dilakukan dalam~$\Zminus$).

\end{document}
